\documentclass[a4paper]{article}

% Paquetes básicos
% Español (México) con XeLaTeX: fontspec para fuentes Unicode, polyglossia
% para la división silábica y los nombres del documento en la variante mexicana.
\usepackage{fontspec}
\usepackage{polyglossia}
\setdefaultlanguage[variant=mexican]{spanish}
% Los nombres chinos van como «pinyin (original)»: xeCJK da a los caracteres
% Han su propia fuente; sin ella XeLaTeX los omite con un aviso.
% FandolSong no cubre algunos caracteres de nombres (U+73FA); WenQuanYi Zen Hei sí.
\usepackage[AutoFallBack=true]{xeCJK}
\setCJKmainfont{FandolSong-Regular.otf}
\setCJKfallbackfamilyfont{\CJKrmdefault}{WenQuanYi Zen Hei}
% polyglossia no traduce los nombres de listings.
\AtBeginDocument{\providecommand{\lstlistingname}{}\renewcommand{\lstlistingname}{Listado}%
  \providecommand{\lstlistlistingname}{}\renewcommand{\lstlistlistingname}{Índice de listados}}
\usepackage{amsmath, amssymb}
\usepackage{graphicx}
\usepackage[margin=2.5cm]{geometry}
\usepackage[most]{tcolorbox}
\usepackage{etoolbox}
\usepackage{listings}
\usepackage{hyperref}
\usepackage{booktabs}
\usepackage{subcaption}
\usepackage{float}

% Cajas de resaltado
\newtcolorbox{knowledgebox}[1]{
    enhanced, colback=blue!5!white, colframe=blue!75!black, colbacktitle=blue!75!black,
    coltitle=white, fonttitle=\bfseries, title={#1}, attach boxed title to top left={yshift=-2mm, xshift=2mm},
    boxrule=1pt, sharp corners
}

\newtcolorbox{importantbox}[1]{
    enhanced, colback=yellow!10!white, colframe=yellow!80!black, colbacktitle=yellow!80!black,
    coltitle=black, fonttitle=\bfseries, title={#1}, sharp corners
}

\newtcolorbox{warningbox}[1]{
    enhanced, colback=red!5!white, colframe=red!75!black, colbacktitle=red!75!black,
    coltitle=white, fonttitle=\bfseries, title={#1}, sharp corners
}

% Listados
\lstset{
    language=Python,
    basicstyle=\ttfamily\small,
    keywordstyle=\color{blue},
    stringstyle=\color{red!60!black},
    commentstyle=\color{green!60!black},
    breaklines=true,
    frame=single,
    numbers=left,
    numberstyle=\tiny\color{gray},
    captionpos=b,
    extendedchars=true
}

% Metadatos de la página inicial
\newcommand{\notetitle}{KAIST CS492D: Introducción a los modelos generativos --- GAN y VAE}
\newcommand{\noteauthors}{Elaboradas a partir de la clase de Minhyuk Sung}
\newcommand{\notedate}{Otoño de 2025}
\newcommand{\videochannel}{Minhyuk Sung (KAIST)}
\newcommand{\videopublishdate}{2025}
\newcommand{\videoduration}{65 minutos}
\newcommand{\videourl}{https://www.youtube.com/watch?v=96XoVM7PYtQ}
\newcommand{\videocoverpath}{cover.jpg}
\newcommand{\repourl}{https://github.com/hqhq1025/ai-course-notes}

% Pie de página con información del repositorio
\usepackage{fancyhdr}
\pagestyle{fancy}
\fancyhf{}
\fancyfoot[C]{\thepage}
\fancyfoot[R]{\footnotesize\color{gray}\href{\repourl}{AI Course Notes}}
\renewcommand{\headrulewidth}{0pt}
\renewcommand{\footrulewidth}{0pt}

\begin{document}

\begin{titlepage}
\centering
{\Large Notas de clase\par}
\vspace{1.2cm}
{\huge\bfseries \notetitle\par}
\vspace{0.8cm}
{\large \noteauthors\par}
\vspace{0.3cm}
{\large \notedate\par}
\vspace{1.2cm}

\ifdefempty{\videocoverpath}{
{\small Al generar las notas, indique la ruta local de la portada del video.\par}
}{
\includegraphics[width=0.82\textwidth,height=0.45\textheight,keepaspectratio]{\videocoverpath}\par
}

\vfill
\begin{tcolorbox}[width=0.9\textwidth, colback=black!2!white, colframe=black!60, sharp corners]
\textbf{Autor o canal del video}: \videochannel\par
\textbf{Fecha de publicación}: \videopublishdate\par
\textbf{Duración del video}: \videoduration\par
\ifdefempty{\videourl}{}{
\textbf{Enlace del video}: \href{\videourl}{\nolinkurl{\videourl}}\par
}
\textbf{Página del curso}: \href{https://diffusion.kaist.ac.kr/}{https://diffusion.kaist.ac.kr/}\par
\textbf{Página del proyecto}: \href{\repourl}{\nolinkurl{\repourl}}\par
\end{tcolorbox}
\end{titlepage}

\tableofcontents
\newpage

%% --- Inicio del contenido --- %%


\section{Resumen del curso y el problema central de los modelos generativos}

Esta es la primera clase del curso ``Diffusion and Flow Models'' de KAIST CS492D. Antes de abordar formalmente los modelos Diffusion y los modelos Flow, el docente Minhyuk Sung repasa los fundamentos de los modelos generativos (Generative Models), cubriendo dos paradigmas clásicos: la teoría de muestreo, las redes generativas adversarias (GANs) y los autoencoders variacionales (VAEs), así como las herramientas de probabilidad y estadística que sustentan a estos modelos.

\subsection{¿Qué son los modelos generativos?}

La tarea central de un modelo generativo puede resumirse en una frase: \textbf{dado un conjunto de muestras de datos, aprender un modelo que sea capaz de generar nuevas muestras que parezcan similares a los datos de entrenamiento pero que no sean idénticas}.

Tomemos como ejemplo las imágenes. Supongamos que disponemos de un conjunto de datos masivo que contiene 58 mil millones de imágenes (por ejemplo, LAION-5B). El objetivo del modelo generativo es entrenar un modelo que pueda producir imágenes completamente nuevas y realistas: no se trata simplemente de copiar una imagen existente del conjunto de datos ni de generar ruido aleatorio, sino de crear muestras nuevas cuya distribución sea consistente con las imágenes reales.

\begin{importantbox}{Perspectiva estadística de los modelos generativos}
Desde el punto de vista estadístico, el problema de los modelos generativos puede formularse de la siguiente manera: supongamos que todos los puntos de datos $x_1, x_2, \ldots, x_N$ se han obtenido mediante muestreo independiente e idénticamente distribuido (i.i.d.) a partir de alguna distribución de probabilidad desconocida $p_{\text{data}}(x)$. Nuestro objetivo es aprender esta distribución para poder generar nuevos puntos de datos a partir de $p_{\text{data}}(x)$.
\end{importantbox}

Para facilitar la comprensión, el docente parte de un escenario más sencillo: supongamos que tenemos un conjunto de puntos en un plano bidimensional que forman cierta forma específica (por ejemplo, dos medias lunas). El problema entonces se convierte en: ¿cómo podemos construir un modelo que sea capaz de generar nuevos puntos de manera que estos parezcan haber sido muestreados de la misma distribución?

\subsection{De lo simple a lo complejo: las imágenes como vectores de alta dimensión}

Una imagen RGB de $256 \times 256$ píxeles puede aplanarse en un vector de $256 \times 256 \times 3 = 196{,}608$ dimensiones. Por lo tanto, un conjunto de imágenes equivale a un conjunto de puntos en un espacio de alta dimensión, y el problema de generación de imágenes es esencialmente equivalente al muestreo en este espacio de alta dimensión.

\begin{knowledgebox}{Representación vectorial de las imágenes}
Cada imagen RGB de $H \times W$ puede considerarse como un punto en el espacio $\mathbb{R}^{H \times W \times 3}$. El conjunto de imágenes constituye una nube de puntos en este espacio de alta dimensión, que define conjuntamente una distribución de datos $p_{\text{data}}(x)$. La tarea del modelo generativo es muestrear nuevos puntos (es decir, nuevas imágenes) a partir de esta distribución.
\end{knowledgebox}

\subsection{Resumen de la sección}

La idea central de los modelos generativos consiste en modelar los datos como muestras de una cierta distribución de probabilidad y generar nuevos datos aprendiendo dicha distribución. Ya se trate de puntos bidimensionales o de imágenes de alta dimensión, la esencia matemática del problema es la misma: muestrear de una distribución desconocida.

\newpage
\section{Fundamentos de la muestreo}

Antes de discutir cómo aprender distribuciones de datos complejas, el instructor repasa una pregunta fundamental: \textbf{si ya conocemos una distribución de probabilidad, ¿cómo se extraen muestras de ella?}

\subsection{Muestreo por transformación inversa (Inverse Transform Sampling)}

\subsubsection{Caso discreto: Muestreo de una distribución Categorical}

Consideremos una distribución categorical con $K$ estados, donde cada estado $k$ tiene probabilidad $p_k$, y $\sum_{k=1}^K p_k = 1$. El método clásico para extraer muestras de esta distribución es el \textbf{muestreo por transformación inversa}:

\begin{enumerate}
    \item Calcular la función de distribución acumulada (CDF): $F(k) = \sum_{i=1}^{k} p_i$
    \item Extraer un valor $u$ de una distribución uniforme $U \sim \text{Uniform}(0, 1)$
    \item Encontrar el estado $k$ que satisface $F(k-1) < u \leq F(k)$; dicho estado es el resultado del muestreo
\end{enumerate}

Intuitivamente, los estados con mayor probabilidad ocupan un ``intervalo'' más ancho en el eje $y$ de la CDF, por lo que la probabilidad de que una muestra uniforme caiga dentro de ese intervalo también es mayor.

\subsubsection{Caso continuo}

Para una variable aleatoria continua $X$, cuya función de densidad de probabilidad (PDF) es $f(x)$ y cuya CDF es:

$$F(x) = \int_{-\infty}^{x} f(t) \, dt$$

Los pasos del muestreo por transformación inversa son:

\begin{enumerate}
    \item Extraer $u \sim \text{Uniform}(0, 1)$
    \item Calcular $x = F^{-1}(u)$
\end{enumerate}

\begin{importantbox}{Fórmula clave del muestreo por transformación inversa}
Sea $F$ la CDF de una distribución continua y $F^{-1}$ su función inversa, entonces:
$$X = F^{-1}(U), \quad U \sim \text{Uniform}(0,1) \implies X \sim F$$
Suposición clave: \textbf{es posible calcular la función inversa de la CDF, $F^{-1}$}.
\end{importantbox}

\subsubsection{Ejercicio de clase: Muestreo de $f(x) = \frac{3}{8}x^2$}

El instructor propone un pequeño ejercicio: dado el PDF $f(x) = \frac{3}{8}x^2$, con $x \in {[}-2, 2{]}$, derive la fórmula de muestreo por transformación inversa.

Solución:
\begin{enumerate}
    \item Calcular la CDF: $F(x) = \int_{-2}^{x} \frac{3}{8}t^2 \, dt = \frac{1}{8}(x^3 + 8)$
    \item Plantear $u = F(x)$ y despejar $x$: $u = \frac{1}{8}(x^3 + 8) \implies x = (8u - 8)^{1/3} = 2(u-1)^{1/3}$
    \item Corrección, derivación más cuidadosa: $x^3 = 8u - 8 \implies x = \sqrt[3]{8u - 8}$
\end{enumerate}

Por lo tanto, el proceso de muestreo consiste en: primero extraer $u \sim \text{Uniform}(0,1)$ y luego calcular $x = (8u - 8)^{1/3}$.

\subsection{Muestreo de la distribución gaussiana: Transformación Box-Muller}

\subsubsection{Planteamiento del problema}

El PDF de la distribución normal estándar $\mathcal{N}(0, 1)$ es:

$$f(x) = \frac{1}{\sqrt{2\pi}} e^{-x^2/2}$$

Su CDF $\Phi(x) = \int_{-\infty}^{x} f(t) \, dt$ \textbf{no tiene una función inversa en forma analítica}, por lo que no se puede utilizar directamente el muestreo por transformación inversa.

\begin{warningbox}{Trampa de la CDF gaussiana irreversible}
Aunque el PDF de la distribución gaussiana tiene una forma analítica elegante, su función inversa (función probit) no existe en forma cerrada. Esto significa que no se puede aplicar directamente el método más simple de muestreo por transformación inversa para distribuciones gaussianas. En la práctica es necesario recurrir a trucos como sustitución de variables.
\end{warningbox}

\subsubsection{Derivación de la Transformación Box-Muller}

La idea central de la Transformación Box-Muller es: \textbf{considerar simultáneamente dos muestras independientes de distribución normal estándar, aplicando una transformación a coordenadas polares para que cada componente pueda muestrearse mediante transformación inversa}.

Sea $X, Y \overset{\text{i.i.d.}}{\sim} \mathcal{N}(0, 1)$, cuya densidad conjunta es:

$$f_{X,Y}(x, y) = \frac{1}{2\pi} e^{-(x^2 + y^2)/2}$$

Introduciendo variables en coordenadas polares $D = X^2 + Y^2$ (radio al cuadrado) y $\Theta$ (ángulo), se puede demostrar:

\begin{itemize}
    \item $\Theta \sim \text{Uniform}(0, 2\pi)$ (distribución uniforme del ángulo, que se extrae directamente)
    \item $D \sim \text{Exponential}(1/2)$, es decir, $f_D(d) = \frac{1}{2}e^{-d/2}$ (distribución $\chi^2$ con 2 grados de libertad)
\end{itemize}

Para la distribución de $D$, la CDF es $F_D(d) = 1 - e^{-d/2}$, cuya función inversa es:

$$d = F_D^{-1}(u) = -2\ln(1 - u)$$

Dado que $1-u$ y $u$ son ambos $\text{Uniform}(0,1)$, esto se simplifica a $d = -2\ln(u)$.

\begin{importantbox}{Fórmula completa de la Transformación Box-Muller}
Dadas dos muestras independientes de distribución uniforme $U_1, U_2 \sim \text{Uniform}(0,1)$:
$$X = \sqrt{-2\ln U_1} \cdot \cos(2\pi U_2)$$
$$Y = \sqrt{-2\ln U_1} \cdot \sin(2\pi U_2)$$
Entonces $X$ y $Y$ son dos muestras independientes de distribución normal estándar. Una sola transformación genera simultáneamente dos muestras gaussianas.
\end{importantbox}

\subsection{Muestreo de distribuciones gaussianas multivariantes y Reparameterization}

\subsubsection{Distribución gaussiana isótropa}

Para una distribución normal multivariante estándar $\mathcal{N}(\mathbf{0}, \mathbf{I})$, dado que las dimensiones son independientes, se puede aplicar la Transformación Box-Muller a cada dimensión por separado.

\subsubsection{Distribución gaussiana anisótropa}

Para una distribución gaussiana multivariante general $\mathcal{N}(\boldsymbol{\mu}, \boldsymbol{\Sigma})$, el método de muestreo es la \textbf{reparametrización}:

$$\mathbf{x} = \boldsymbol{\mu} + \mathbf{L} \boldsymbol{\epsilon}, \quad \boldsymbol{\epsilon} \sim \mathcal{N}(\mathbf{0}, \mathbf{I})$$

Donde $\mathbf{L}$ satisface $\boldsymbol{\Sigma} = \mathbf{L}\mathbf{L}^T$ (por ejemplo, descomposición de Cholesky).

\begin{importantbox}{Truco de reparametrización}
Muestrear de $\mathcal{N}(\boldsymbol{\mu}, \boldsymbol{\Sigma})$ es equivalente a:
$$\mathbf{x} = \boldsymbol{\mu} + \boldsymbol{\Sigma}^{1/2} \boldsymbol{\epsilon}, \quad \boldsymbol{\epsilon} \sim \mathcal{N}(\mathbf{0}, \mathbf{I})$$
Este truco no solo es fundamental para el muestreo, sino clave para implementar la propagación del gradiente en VAEs: aísla la aleatoriedad en $\boldsymbol{\epsilon}$, haciendo que el proceso de muestreo sea diferenciable respecto a $\boldsymbol{\mu}$ y $\boldsymbol{\Sigma}$.
\end{importantbox}

\begin{knowledgebox}{¿Por qué los modelos generadores prefieren las distribuciones gaussianas?}
Las distribuciones gaussianas son omnipresentes en los modelos generativos debido a: (1) conocemos cómo muestrearlas de manera eficiente; (2) poseen excelentes propiedades matemáticas (como que la divergencia KL tiene forma cerrada); (3) el teorema del límite central garantiza que muchas distribuciones tienden asintóticamente a la gaussiana; y (4) bajo restricciones de media y varianza, la distribución gaussiana maximiza la entropía.
\end{knowledgebox}

\subsection{Resumen de la sección}

Cuando la forma de una distribución de probabilidad es conocida, podemos realizar muestreo mediante transformación inversa, sustitución de variables (como la Transformación Box-Muller) y reparametrización. Sin embargo, las distribuciones de datos reales (como la de imágenes naturales) tienen formas desconocidas y extremadamente complejas; esto es precisamente lo que motiva el uso de modelos generativos profundos.

\newpage
\section{Marco unificado para modelos generativos: de Latent a Data}

\subsection{Idea central: aprender mapas entre distribuciones}

\begin{importantbox}{Enfoque común en todos los modelos generativos}
Todos los modelos generativos profundos (GAN, VAE, Diffusion Model, Flow Model) comparten la misma idea central: \textbf{aprender un mapa desde una distribución simple (latent distribution) hacia una distribución de datos compleja (data distribution)}.

En particular:
\begin{enumerate}
    \item Seleccionar una distribución simple fácil de muestrear como latent distribution $p(z)$ (generalmente una distribución gaussiana estándar $\mathcal{N}(\mathbf{0}, \mathbf{I})$)
    \item Aprender un mapa mediante una red neuronal $G_\theta: z \mapsto x$
    \item Durante la generación: muestrear $z \sim p(z)$, calcular $x = G_\theta(z)$
\end{enumerate}
La diferencia entre los distintos modelos radica en \textbf{cómo entrenar este mapa}.
\end{importantbox}

\subsection{Convenciones de notación}

Este curso utiliza las siguientes notaciones que se mantienen a lo largo del texto:

\begin{table}[H]
\centering
\begin{tabular}{lll}
\toprule
Símbolo & Significado & Comentario \\
\midrule
$z$ & variable latent & proviene de una distribución simple (generalmente $\mathcal{N}(\mathbf{0}, \mathbf{I})$) \\
$x$ & punto de datos & proviene del conjunto de datos de entrenamiento \\
$p(z)$ & distribución latent / prior & conocida \\
$p(x)$ / $p_{\text{data}}(x)$ & distribución de datos & \textbf{desconocida}; este es el objetivo que queremos aproximar \\
$p_\theta(x|z)$ & distribución likelihood & distribución condicional para generar $x$ dado $z$ \\
$p(z|x)$ & distribución posterior & distribución condicional para inferir $z$ dado $x$ \\
\bottomrule
\end{tabular}
\caption{Convenciones de notación utilizadas en este curso}
\end{table}

\subsection{La inspiración del Autoencoder}

El autoencoder es un punto de partida para comprender los modelos generativos:

\begin{itemize}
    \item \textbf{Encoder}: $f_\phi: x \mapsto z$, comprime el punto de datos hacia un espacio latent de baja dimensión
    \item \textbf{Decoder}: $g_\theta: z \mapsto \hat{x}$, reconstruye los datos a partir de la representación latent
    \item Objetivo del entrenamiento: minimizar el error de reconstrucción $\|x - g_\theta(f_\phi(x))\|^2$
\end{itemize}

En el contexto de los modelos generativos, el \textbf{decoder es el mapa que necesitamos}; mapea puntos del espacio latent hacia el espacio de datos. Pero la pregunta clave es: \textbf{cómo garantizar que el decoder pueda mapear cualquier punto extraído aleatoriamente de la latent distribution en un punto de datos razonable}?

\begin{warningbox}{Un autoencoder no es un modelo generativo}
El espacio latent de un autoencoder estándar no sigue alguna distribución simple conocida. Si tomamos directamente un punto $z$ al azar del espacio latent y lo introducimos en el decoder, la salida resultante probablemente no parecerá datos reales. VAE resuelve este problema imponiendo restricciones previas al espacio latent.
\end{warningbox}

\subsection{Resumen de la sección}

El paradigma unificado de los modelos generativos es: partiendo de una distribución simple conocida $p(z)$, aprender mediante redes neuronales el mapa hacia la distribución de datos $p_{\text{data}}(x)$. La única diferencia entre GAN, VAE, Diffusion Model y Flow Model radica en sus estrategias de entrenamiento. El decoder del autoencoder proporciona un prototipo de este mapa, pero requiere mecanismos adicionales para garantizar la estructura del espacio latent.

\newpage
\section{GAN: Redes generativas adversarias}

\subsection{Del ataque adversario a los modelos generativos}

La idea detrás de las GANs (Generative Adversarial Networks) se origina en el ataque adversario. El docente ilustró esto con un ejemplo clásico: un clasificador entrenado puede identificar correctamente un panda, pero si se añade una perturbación microscópica casi imperceptible para el ojo humano (ruido) a la imagen, el clasificador la clasificará erróneamente como un avestruz —incluso si se amplía la diferencia 10 veces, es difícil que el ojo humano distinga la diferencia.

\begin{knowledgebox}{Del ataque adversario al salto conceptual de las GANs}
El ataque adversario reveló una verdad: las redes neuronales pueden ser extremadamente sensibles a cambios sutiles imperceptibles para el ojo humano. Ian Goodfellow, el creador de las GANs, se inspiró en esto: si una red (atacante/generador) puede generar muestras que engañen a otra red (defensor/discriminador), entonces mediante la competencia entre ambas, el generador eventualmente aprenderá a producir datos realistas.
\end{knowledgebox}

\subsection{Arquitectura de las GANs}

Las GANs están compuestas por dos redes neuronales:

\begin{itemize}
    \item \textbf{Generator (Generador)} $G_\theta$: recibe muestras latentes $z \sim p(z)$ y produce puntos de datos ``falsos'' $G_\theta(z)$.
    \item \textbf{Discriminator (Discriminador)} $D_\phi$: recibe un punto de datos y devuelve la probabilidad de que ese punto sea ``datos reales'', $D_\phi(x) \in {[}0, 1{]}$.
\end{itemize}

El generador y el discriminador compiten entre sí:
\begin{itemize}
    \item Objetivo del discriminador: distinguir correctamente los datos reales de los generados.
    \item Objetivo del generador: producir muestras que engañen al discriminador.
\end{itemize}

\subsection{Función de pérdida de las GANs: Juego Minimax}

El objetivo de entrenamiento de las GANs es un problema de optimización minimax:

$$\min_\theta \max_\phi \; V(G_\theta, D_\phi) = \mathbb{E}_{x \sim p_{\text{data}}} {[}\log D_\phi(x){]} + \mathbb{E}_{z \sim p(z)} {[}\log(1 - D_\phi(G_\theta(z))){]}$$

\begin{importantbox}{Significado de cada término en la función de pérdida de las GANs}
\begin{itemize}
    \item $\mathbb{E}_{x \sim p_{\text{data}}} {[}\log D_\phi(x){]}$: para los datos reales, el discriminador debe devolver una probabilidad alta (cercana a 1), haciendo que $\log D_\phi(x)$ sea lo más grande posible.
    \item $\mathbb{E}_{z \sim p(z)} {[}\log(1 - D_\phi(G_\theta(z))){]}$: para los datos generados $G_\theta(z)$, el discriminador debe devolver una probabilidad baja (cercana a 0), haciendo que $1 - D_\phi(G_\theta(z))$ se acerque a 1.
\end{itemize}
El \textbf{discriminador} mejora su capacidad de distinción \textbf{maximizando} $V$; el \textbf{generador} mejora su capacidad de engaño \textbf{minimizando} $V$.
\end{importantbox}

\subsubsection{Solución óptima del discriminador}

Manteniendo fijo el generador $G_\theta$, buscamos la solución óptima para el discriminador:

$$D^*_\phi(x) = \frac{p_{\text{data}}(x)}{p_{\text{data}}(x) + p_G(x)}$$

Donde $p_G(x)$ es la distribución inducida por el generador. Cuando $p_G = p_{\text{data}}$, entonces $D^*(x) = 1/2$, lo que significa que el discriminador no puede distinguir entre verdadero y falso —esto es precisamente el estado ideal hacia el cual converge el entrenamiento de las GANs.

\subsubsection{GANs con solución óptima global}

Se puede demostrar que cuando $p_G = p_{\text{data}}$, la función de pérdida alcanza su valor mínimo global $-2\log 2$. En este punto, la distribución del generador coincide completamente con la distribución de los datos.

\subsection{Dificultades en el entrenamiento de las GANs}

\begin{warningbox}{Problemas comunes en el entrenamiento de las GANs}
El problema de optimización minimax es \textbf{extremadamente difícil de optimizar} en la práctica y tiende a caer en óptimos locales pobres. Los modos de falla más frecuentes incluyen:
\begin{itemize}
    \item \textbf{Discriminador dominante (Discriminator Dominance)}: el discriminador converge mucho más rápido que el generador, lo que resulta en una señal de gradiente extremadamente débil (casi nula) para el generador, impidiendo un aprendizaje efectivo. Este es el modo de falla más común.
    \item \textbf{Colapso de modos (Mode Collapse)}: incluso si el entrenamiento progresa bien, el generador podría aprender a producir solo una parte de los modos de la distribución de datos, perdiendo así la diversidad de los datos. Por ejemplo, en datos multimodales, el generador podría concentrarse produciendo muestras de un solo modo.
    \item \textbf{Inestabilidad en el entrenamiento}: la curva de pérdida oscila violentamente, haciendo difícil determinar si se ha alcanzado la convergencia.
\end{itemize}
\end{warningbox}

El docente mencionó una analogía vívida: si estableces criterios de éxito extremadamente estrictos desde el principio (un discriminador demasiado fuerte), no aprenderás nada —esto es exactamente como las lecciones de la vida.

\subsubsection{Práctica ingenieril en el entrenamiento de las GANs}

En la práctica, entrenar GANs requiere numerosas técnicas de ingeniería:

\begin{itemize}
    \item Por lo general, se preentrena al generador durante unos pasos para que pueda producir salidas con algún significado antes de introducir el discriminador para el entrenamiento adversarial.
    \item Se necesita equilibrar cuidadosamente las tasas de aprendizaje y la proporción de pasos de entrenamiento de ambas redes.
    \item Muchas empresas dedicaron grandes equipos de ingeniería especializados a ajustar los parámetros de entrenamiento de las GANs antes del año 2020.
\end{itemize}

\subsection{Evolución de las GANs}

Desde que Goodfellow y sus colegas propusieron las GANs en 2013/2014, han surgido numerosas variantes. El docente mencionó brevemente varios hitos fundamentales:

\begin{table}[H]
\centering
\begin{tabular}{lll}
\toprule
Modelo & Año & Contribución principal \\
\midrule
GAN & 2014 & Propuesta del marco de entrenamiento adversario \\
DCGAN & 2015 & Introducción de la arquitectura convolucional \\
WGAN & 2017 & Sustitución de la divergencia JS por la distancia Wasserstein \\
Progressive GAN & 2017 & Estrategia de crecimiento progresivo \\
StyleGAN / StyleGAN2 & 2019/2020 & Generación de rostros de alta calidad \\
\bottomrule
\end{tabular}
\caption{Variantes importantes de las GANs (parcial)}
\end{table}

\begin{knowledgebox}{¿Por qué se pasó de las GANs a los modelos de difusión?}
El docente señaló que la razón principal por la cual la industria pasó de las GANs a los modelos de Difusión/Flujo no fue la calidad —StyleGAN ya podía generar imágenes de rostros extremadamente realistas en 2019—. El verdadero impulso fue: (1) \textbf{Dificultad en el entrenamiento} —la optimización minimax es inestable y requiere una gran inversión de ingeniería; (2) \textbf{Falta de diversidad} —el problema del colapso de modos hace que las muestras generadas carezcan de variedad. El modelo de difusión evita astutamente estos dos problemas al descomponer el proceso de generación en múltiples pasos de eliminación de ruido.
\end{knowledgebox}

\subsection{Resumen de la sección}

Las GANs aprenden la distribución de datos mediante un juego adversario entre un generador y un discriminador. Su marco teórico es elegante (el juego minimax puede coincidir exactamente con la distribución de datos en condiciones ideales), pero en la práctica enfrenta problemas graves como la inestabilidad del entrenamiento y el colapso de modos. Estos desafíos impulsaron el desarrollo de alternativas como las VAEs y, más tarde, los modelos de difusión.

\newpage
\section{VAE: variational autoencoder}

\subsection{Motivación de pasar de GAN a VAE}

La dificultad central de los GAN radica en la optimización minimax. La idea detrás del VAE es: \textbf{¿podemos evitar el entrenamiento adversarial y aprender la mapeo de la distribución latente a la distribución de datos mediante una forma más estable?}

La innovación clave del VAE consiste en dejar de aprender una función de mapeo determinista $G_\theta(z)$ para aprender una \textbf{distribución de probabilidad condicional} $p_\theta(x|z)$, es decir, la distribución de los datos $x$ dado que tenemos la variable latente $z$.

\subsection{Marco probabilístico del VAE}

\subsubsection{Proceso generativo}

El VAE define el siguiente proceso generativo:

\begin{enumerate}
    \item Muestreo desde una distribución a priori: $z \sim p(z) = \mathcal{N}(\mathbf{0}, \mathbf{I})$
    \item Muestreo desde la distribución de verosimilitud: $x \sim p_\theta(x|z) = \mathcal{N}(\mu_\theta(z), \sigma^2 \mathbf{I})$
\end{enumerate}

Donde $\mu_\theta(z)$ es la salida de la red neuronal del decodificador (la predicción del valor medio de $x$ dado $z$) y $\sigma^2$ es un hiperparámetro de varianza fija.

\begin{importantbox}{Modelo probabilístico generativo del VAE}
$$p_\theta(x, z) = p_\theta(x|z) \cdot p(z)$$
Donde:
\begin{itemize}
    \item $p(z) = \mathcal{N}(\mathbf{0}, \mathbf{I})$: distribución a priori (prior), conocida.
    \item $p_\theta(x|z) = \mathcal{N}(\mu_\theta(z), \sigma^2 \mathbf{I})$: distribución de verosimilitud (likelihood), parametrizada por el decodificador.
    \item $p_\theta(x) = \int p_\theta(x|z) p(z) \, dz$: verosimilitud marginal (marginal likelihood), obtenida integrando respecto a $z$.
\end{itemize}
\end{importantbox}

\subsubsection{Objetivo de entrenamiento: maximizar la verosimilitud marginal}

En principio, deseamos maximizar el logaritmo de la verosimilitud marginal de los datos:

$$\max_\theta \; \mathbb{E}_{x \sim p_{\text{data}}} [\log p_\theta(x)] = \max_\theta \; \mathbb{E}_{x \sim p_{\text{data}}} \left[\log \int p_\theta(x|z) p(z) \, dz \right]$$

\begin{warningbox}{La verosimilitud marginal no es calculable}
La integral $p_\theta(x) = \int p_\theta(x|z) p(z) \, dz$ \textbf{no puede calcularse de forma exacta} en espacios de alta dimensión. Intuitivamente: la dimensión del espacio latente podría ser de cientos de dimensiones, e integrar sobre todos los posibles $z$ no es factible. Por eso necesitamos inferencia variacional (Variational Inference).
\end{warningbox}

\subsection{Repaso de probabilidad: herramientas matemáticas necesarias para el VAE}

\subsubsection{Marginalización}

Dada una distribución conjunta $p(x, z)$, se puede obtener la distribución marginal de otra variable integrando (o sumando) respecto a una de ellas:

$$p(x) = \int p(x, z) \, dz = \int p(x|z) p(z) \, dz$$

\subsubsection{Teorema de Bayes}

$$p(z|x) = \frac{p(x|z) p(z)}{p(x)}$$

Convención de terminología para cada término:

\begin{table}[H]
\centering
\begin{tabular}{ll}
\toprule
Símbolo & Nombre \\
\midrule
$p(z)$ & Prior (a priori): conocimiento a priori sobre $z$ \\
$p(x|z)$ & Likelihood (verosimilitud): probabilidad de observar $x$ dado $z$ \\
$p(z|x)$ & Posterior (a posteriori): creencia actualizada sobre $z$ tras observar $x$ \\
$p(x)$ & Evidence / Marginal Likelihood (evidencia): probabilidad marginal de los datos \\
\bottomrule
\end{tabular}
\caption{Terminología de los términos en el teorema de Bayes}
\end{table}

\subsubsection{Divergencia de Kullback-Leibler}

La divergencia de Kullback-Leibler mide la "distancia" entre dos distribuciones de probabilidad (nota: no es simétrica, por lo que no es estrictamente una métrica de distancia):

$$D_{\mathrm{KL}}(p \| q) = \mathbb{E}_{x \sim p} \left[\log \frac{p(x)}{q(x)}\right] = \int p(x) \log \frac{p(x)}{q(x)} \, dx$$

\begin{importantbox}{Propiedades clave de la divergencia de KL}
\begin{itemize}
    \item \textbf{No negatividad}: $D_{\mathrm{KL}}(p \| q) \geq 0$, con igualdad si y solo si $p = q$ casi en todas partes.
    \item \textbf{Asimetría}: generalmente $D_{\mathrm{KL}}(p \| q) \neq D_{\mathrm{KL}}(q \| p)$.
    \item \textbf{Forma cerrada para la divergencia KL entre dos distribuciones gaussianas}.
\end{itemize}
\end{importantbox}

\begin{knowledgebox}{Divergencia KL entre dos distribuciones gaussianas multivariadas}
Supongamos $p = \mathcal{N}(\boldsymbol{\mu}_1, \boldsymbol{\Sigma}_1)$ y $q = \mathcal{N}(\boldsymbol{\mu}_2, \boldsymbol{\Sigma}_2)$. Entonces:
$$D_{\mathrm{KL}}(p \| q) = \frac{1}{2}\left[\log\frac{|\boldsymbol{\Sigma}_2|}{|\boldsymbol{\Sigma}_1|} - d + \mathrm{tr}(\boldsymbol{\Sigma}_2^{-1}\boldsymbol{\Sigma}_1) + (\boldsymbol{\mu}_2 - \boldsymbol{\mu}_1)^T \boldsymbol{\Sigma}_2^{-1} (\boldsymbol{\mu}_2 - \boldsymbol{\mu}_1)\right]$$
donde $d$ es la dimensión. En particular, cuando $q = \mathcal{N}(\mathbf{0}, \mathbf{I})$:
$$D_{\mathrm{KL}}(\mathcal{N}(\boldsymbol{\mu}, \boldsymbol{\Sigma}) \| \mathcal{N}(\mathbf{0}, \mathbf{I})) = \frac{1}{2}\left[-\log|\boldsymbol{\Sigma}| - d + \mathrm{tr}(\boldsymbol{\Sigma}) + \boldsymbol{\mu}^T\boldsymbol{\mu}\right]$$
Esta fórmula se usa directamente en el entrenamiento del VAE.
\end{knowledgebox}

\subsection{Derivación de la desigualdad de Jensen y el ELBO}

\subsubsection{Funciones convexas e inecuación de Jensen}

\textbf{Función convexa (Convex Function)}: se define como aquella que, para cualquier par de puntos $x_1, \ldots, x_n$ y pesos $t_1, \ldots, t_n$ ($t_i \geq 0$, $\sum t_i = 1$), cumple:

$$f\left(\sum_i t_i x_i\right) \leq \sum_i t_i f(x_i)$$

Intuitivamente, el gráfico de una función convexa se curva "hacia abajo", y cualquier segmento que une dos puntos yace sobre el gráfico de la función.

Una \textbf{función cóncava (Concave Function)} es lo opuesto a la convexa: su gráfico se curva "hacia arriba" y $-f$ es una función convexa. La función $\log$ es una función cóncava importante.

\begin{importantbox}{Inecuación de Jensen}
Si $f$ es una función convexa y $X$ es una variable aleatoria, entonces:
$$f(\mathbb{E}{[}X{]}) \leq \mathbb{E}{[}f(X){]}$$
Si $f$ es cóncava (como $\log$), la dirección de la inecuación se invierte:
$$f(\mathbb{E}{[}X{]}) \geq \mathbb{E}{[}f(X){]}$$
Es decir, $\log(\mathbb{E}{[}X{]}) \geq \mathbb{E}{[}\log X{]}$.
\end{importantbox}

\subsubsection{Derivación del ELBO}

El ELBO (Evidence Lower Bound) es el objetivo central del entrenamiento del VAE. A continuación se muestra el proceso de derivación completo.

\textbf{Punto de partida}: queremos maximizar el logaritmo de la verosimilitud marginal $\log p_\theta(x)$, pero no puede calcularse directamente. Introduciendo una distribución a posteriori aproximada $q_\phi(z|x)$ (encoder), tenemos:

\begin{align}
\log p_\theta(x) &= \log \int p_\theta(x, z) \, dz \nonumber \\
&= \log \int \frac{p_\theta(x, z)}{q_\phi(z|x)} q_\phi(z|x) \, dz \nonumber \\
&= \log \mathbb{E}_{z \sim q_\phi(z|x)} \left[\frac{p_\theta(x, z)}{q_\phi(z|x)}\right] \nonumber \\
&\geq \mathbb{E}_{z \sim q_\phi(z|x)} \left[\log \frac{p_\theta(x, z)}{q_\phi(z|x)}\right] \quad \text{(Inecuación de Jensen, $\log$ es cóncava)}
\end{align}

La última línea es el ELBO:

$$\text{ELBO}(\theta, \phi; x) = \mathbb{E}_{z \sim q_\phi(z|x)} \left[\log \frac{p_\theta(x, z)}{q_\phi(z|x)}\right]$$

\begin{importantbox}{Descomposición del ELBO}
Sustituyendo $p_\theta(x, z) = p_\theta(x|z)p(z)$, el ELBO se puede descomponer en dos términos:
\begin{align}
\text{ELBO} &= \mathbb{E}_{z \sim q_\phi(z|x)} \left[\log p_\theta(x|z)\right] - D_{\mathrm{KL}}(q_\phi(z|x) \| p(z)) \nonumber
\end{align}
\begin{itemize}
    \item \textbf{Término de reconstrucción (Reconstruction Term)}: $\mathbb{E}_{z \sim q_\phi(z|x)} {[}\log p_\theta(x|z){]}$ —mide la capacidad del decodificador para reconstruir $x$ a partir de $z$. Cuando $p_\theta(x|z) = \mathcal{N}(\mu_\theta(z), \sigma^2 I)$, este término es equivalente al error cuadrático medio (MSE) negativo.
    \item \textbf{Término de regularización (Regularization / KL Term)}: $-D_{\mathrm{KL}}(q_\phi(z|x) \| p(z))$ —incentiva que la distribución de salida del encoder se acerque a la a priori $p(z) = \mathcal{N}(\mathbf{0}, \mathbf{I})$, estructurando así el espacio latente.
\end{itemize}
\end{importantbox}

\subsubsection{Relación entre ELBO y verosimilitud marginal}

En realidad, se puede escribir con más precisión:

$$\log p_\theta(x) = \text{ELBO}(\theta, \phi; x) + D_{\mathrm{KL}}(q_\phi(z|x) \| p_\theta(z|x))$$

\begin{warningbox}{El ELBO es una cota inferior, no una igualdad}
Dado que $D_{\mathrm{KL}}(q_\phi(z|x) \| p_\theta(z|x)) \geq 0$, el ELBO siempre será menor o igual al logaritmo de la verosimilitud marginal real $\log p_\theta(x)$. La diferencia (gap) es igual a la divergencia KL entre la a posteriori aproximada $q_\phi(z|x)$ y la a posteriori real $p_\theta(z|x)$. Cuando $q_\phi(z|x) = p_\theta(z|x)$ (la a posteriori aproximada coincide exactamente con la real), el ELBO es igual a $\log p_\theta(x)$ y el gap es cero.
\end{warningbox}

\subsection{Arquitectura de la red del VAE}

En términos de implementación, el VAE incluye:

\begin{itemize}
    \item \textbf{Red Encoder} $q_\phi(z|x)$: entrada $x$, salida $\mu_\phi(x)$ y $\sigma_\phi(x)$ (media y desviación estándar de la distribución latente).
    \item \textbf{Red Decoder} $p_\theta(x|z)$: entrada $z$, salida $\mu_\theta(z)$ (valor medio para reconstruir $x$).
    \item \textbf{Proceso de muestreo}: se realiza mediante el truco de reparametrización para obtener $z = \mu_\phi(x) + \sigma_\phi(x) \odot \epsilon$, con $\epsilon \sim \mathcal{N}(\mathbf{0}, \mathbf{I})$.
\end{itemize}

\begin{knowledgebox}{Papel clave del truco de reparametrización en el VAE}
La operación de muestrear $z$ directamente desde $q_\phi(z|x) = \mathcal{N}(\mu_\phi(x), \sigma^2_\phi(x) \mathbf{I})$ no es diferenciable y no permite optimizar $\phi$ mediante retropropagación. El truco de reparametrización expresa el muestreo como $z = \mu_\phi(x) + \sigma_\phi(x) \odot \epsilon$ (donde $\epsilon \sim \mathcal{N}(\mathbf{0}, \mathbf{I})$ es independiente de los parámetros), permitiendo que el gradiente fluya a través de $\mu_\phi$ y $\sigma_\phi$, logrando así un entrenamiento end-to-end.
\end{knowledgebox}

\subsection{Función de pérdida del VAE}

Sintetizando las derivaciones anteriores, el objetivo del entrenamiento del VAE es maximizar el ELBO, lo cual equivale a minimizar la siguiente función de pérdida:

$$\mathcal{L}_{\text{VAE}}(\theta, \phi; x) = -\text{ELBO} = -\mathbb{E}_{z \sim q_\phi(z|x)}{[}\log p_\theta(x|z){]} + D_{\mathrm{KL}}(q_\phi(z|x) \| p(z))$$

Cuando $p_\theta(x|z) = \mathcal{N}(\mu_\theta(z), \sigma^2 \mathbf{I})$ y $p(z) = \mathcal{N}(\mathbf{0}, \mathbf{I})$:

$$\mathcal{L}_{\text{VAE}} = \frac{1}{2\sigma^2}\|x - \mu_\theta(z)\|^2 + D_{\mathrm{KL}}(\mathcal{N}(\mu_\phi(x), \text{diag}(\sigma^2_\phi(x))) \| \mathcal{N}(\mathbf{0}, \mathbf{I}))$$

\begin{importantbox}{Intuición detrás de la función de pérdida del VAE}
La pérdida del VAE consta de dos partes:
\begin{enumerate}
    \item \textbf{Error de reconstrucción (MSE)}: obliga al modelo a reconstruir con precisión los datos de entrada.
    \item \textbf{Penalización KL}: fuerza a que la distribución latente se acerque a una gaussiana estándar, haciendo que el espacio latente sea suave, continuo y muestreable.
\end{enumerate}
Existe tensión entre ambos términos: el error de reconstrucción desea que las codificaciones latentes conserven al máximo la información (lo que podría llevar a distribuciones latentes superpuestas para diferentes puntos de datos), mientras que el término KL empuja a todas las distribuciones latentes hacia la gaussiana estándar (lo que podría provocar pérdida de información). Encontrar un buen equilibrio es el núcleo del ajuste de hiperparámetros en el VAE.
\end{importantbox}

\subsection{Relación entre el VAE y los modelos de difusión}

Durante la explicación del VAE, el profesor de clase adelantó varias veces: \textbf{la idea del VAE está directamente relacionada con los Modelos de Difusión}. En concreto:

\begin{itemize}
    \item Los Modelos de Difusión pueden entenderse como un tipo de VAE "jerárquico", donde las variables latentes son una serie de estados intermedios con ruido progresivo.
    \item El método de derivación del ELBO puede generalizarse directamente a los Modelos de Difusión.
    \item El truco de reparametrización presente en el VAE es igualmente indispensable en los Modelos de Difusión.
\end{itemize}

\subsection{Resumen de la sección}

El VAE evita las dificultades de optimización minimax de los GAN mediante inferencia variacional. Introduce un encoder (a posteriori aproximada) y un decodificador (verosimilitud), entrenando maximizando el ELBO (cota inferior del logaritmo de la verosimilitud marginal). El ELBO se descompone naturalmente en un término de reconstrucción y un término de regularización KL; el equilibrio entre ambos determina tanto la calidad generativa como el grado de estructuración del espacio latente. El entrenamiento del VAE es considerablemente más estable que el de los GAN, y su marco probabilístico sienta las bases teóricas para los posteriores Modelos de Difusión.

\newpage
\section{Caja de herramientas de probabilidad}

Esta sección resume los conceptos y fórmulas clave de probabilidad tratados en esta clase, para consulta rápida.

\subsection{Marginalización}

$$p(x) = \int p(x, z) \, dz = \int p(x|z) p(z) \, dz$$

En el caso discreto, la integral se reemplaza por una suma: $p(x) = \sum_z p(x, z)$

\subsection{Probabilidad condicional y teorema de Bayes}

$$p(z|x) = \frac{p(x|z) p(z)}{p(x)} = \frac{p(x|z) p(z)}{\int p(x|z') p(z') \, dz'}$$

\subsection{Esperanza}

$$\mathbb{E}_{x \sim p}{[}f(x){]} = \int f(x) p(x) \, dx$$

Para variables aleatorias discretas: $\mathbb{E}{[}X{]} = \sum_i x_i p(x_i)$

\subsection{Divergencia de Kullback-Leibler}

$$D_{\mathrm{KL}}(p \| q) = \int p(x) \log \frac{p(x)}{q(x)} \, dx \geq 0$$

La igualdad se cumple si y solo si $p = q$ (casi en todas partes).

\subsection{Desigualdad de Jensen}

Para funciones convexas $f$: $f(\mathbb{E}{[}X{]}) \leq \mathbb{E}{[}f(X){]}$

Para funciones cóncavas $f$ (como $\log$): $f(\mathbb{E}{[}X{]}) \geq \mathbb{E}{[}f(X){]}$

\subsection{Resumen de la sección}

La marginalización, el teorema de Bayes, la divergencia de Kullback-Leibler y la desigualdad de Jensen son los fundamentos matemáticos para entender VAE y Diffusion Model. En particular, la desigualdad de Jensen deriva directamente el ELBO, mientras que la divergencia de Kullback-Leibler aparece tanto en la descomposición del ELBO como en el término de regularización de VAE.

\newpage
\section{Síntesis y ampliación}

\subsection{Puntos clave de esta clase}

Esta clase es la introducción al curso KAIST CS492D, estableciendo sistemáticamente el marco teórico para los modelos generativos:

\begin{enumerate}
    \item \textbf{Problema central de los modelos generativos}: Dado un conjunto de muestras de datos, aprender la distribución subyacente $p_{\text{data}}(x)$ para generar nuevas muestras.
    \item \textbf{Marco unificado}: Todos los modelos de generación profunda aprenden una mapeo desde una distribución simple $p(z)$ hacia la distribución de datos $p_{\text{data}}(x)$.
    \item \textbf{GAN}: Aprende el mapeo mediante un juego entre generador y discriminador (optimización minimax). Teóricamente elegante, pero su entrenamiento es extremadamente inestable.
    \item \textbf{VAE}: Aprende la distribución condicional $p_\theta(x|z)$ mediante inferencia variacional, maximizando el ELBO. Su entrenamiento es estable y sienta las bases para los Diffusion Models.
    \item \textbf{Herramientas matemáticas}: Muestreo por transformación inversa, transformación Box-Muller, truco de reparametrización, marginalización, teorema de Bayes, divergencia KL, desigualdad de Jensen.
\end{enumerate}

\subsection{Comparación GAN vs. VAE}

\begin{table}[H]
\centering
\begin{tabular}{lll}
\toprule
Característica & GAN & VAE \\
\midrule
Método de entrenamiento & Entrenamiento adversario Minimax & Maximización del ELBO \\
Estabilidad del entrenamiento & Inestable, propenso a fallos & Relativamente estable \\
Calidad generada & Alta (especialmente nitidez en imágenes) & Media (tiende a la difuminación) \\
Diversidad & Propenso al colapso de modos & Mejor \\
Fundamento teórico & Teoría de juegos & Inferencia variacional \\
¿Posee encoder? & No (solo generador + discriminador) & Sí (encoder + decoder) \\
Estructura del espacio latente & Sin restricciones explícitas & Con restricción regularizadora KL \\
Relación con Diffusion & Débil & Fuerte (el marco ELBO se generaliza directamente) \\
\bottomrule
\end{tabular}
\caption{Comparación sistemática entre GAN y VAE}
\end{table}

\subsection{Perspectivas: Modelos de Difusión y Flujo}

Esta clase prepara el terreno con toda la base teórica necesaria para los contenidos posteriores. En las siguientes clases, el curso explorará en profundidad:

\begin{itemize}
    \item \textbf{Diffusion Model}: Puede entenderse como un ``VAE de múltiples pasos'', que modela la distribución de datos mediante un proceso gradual de adición y eliminación de ruido. La derivación del ELBO se puede generalizar directamente al contexto de difusión.
    \item \textbf{Flow Model}: Modela la verosimilitud logarítmica exacta mediante transformaciones invertibles (sin necesidad de una aproximación inferior como el ELBO).
    \item \textbf{Score-based Model}: Logra la generación desde otra perspectiva, aprendiendo la función score de la distribución de datos (el gradiente de la densidad logarítmica).
\end{itemize}

Los conceptos teóricos de muestreo de esta clase (muestreo por transformación inversa, truco de reparametrización), el marco probabilístico (prior / verosimilitud / posterior) y la derivación del ELBO aparecerán y se profundizarán en cada una de las siguientes clases.

\subsection{Lecturas adicionales}

\begin{itemize}
    \item Goodfellow et al., \textit{Generative Adversarial Nets}, NeurIPS 2014
    \item Kingma \& Welling, \textit{Auto-Encoding Variational Bayes}, ICLR 2014
    \item Kingma \& Welling, \textit{An Introduction to Variational Autoencoders}, Foundations and Trends in ML, 2019
    \item Doersch, \textit{Tutorial on Variational Autoencoders}, arXiv:1606.05908
    \item Página principal del curso KAIST CS492D: \url{https://diffusion.kaist.ac.kr/}
\end{itemize}

%% --- Fin del contenido --- %%

\end{document}
